\documentclass[10pt,a4paper]{amsart}
\usepackage{amsmath,amssymb,amsthm,graphicx}
% Retained minimal preamble for the Stacks Project chapter sites-cohomology.tex (Cohomology on Sites).
% Verbatim excerpt of the chapter preamble: Xy-pic setup (native preamble lines 14-18) and the native
% theorem declarations with the mathematical macros (native preamble lines 154-217). Copied mechanically;
% nothing added, removed or redefined. The wrapper class and the amsmath/amssymb/amsthm/graphicx base
% packages are supplied by the assembler, so the class/externaldocument part of the native preamble is omitted.
% For commutative diagrams we use Xy-pic
\usepackage[all]{xy}

% We use 2cell for 2-commutative diagrams.
\xyoption{2cell}
% Theorem environments.
%
\theoremstyle{plain}
\newtheorem{theorem}[subsection]{Theorem}
\newtheorem{proposition}[subsection]{Proposition}
\newtheorem{lemma}[subsection]{Lemma}

\theoremstyle{definition}
\newtheorem{definition}[subsection]{Definition}
\newtheorem{example}[subsection]{Example}
\newtheorem{exercise}[subsection]{Exercise}
\newtheorem{situation}[subsection]{Situation}

\theoremstyle{remark}
\newtheorem{remark}[subsection]{Remark}
\newtheorem{remarks}[subsection]{Remarks}

\numberwithin{equation}{subsection}

% Macros
%
\def\lim{\mathop{\mathrm{lim}}\nolimits}
\def\colim{\mathop{\mathrm{colim}}\nolimits}
\def\Spec{\mathop{\mathrm{Spec}}}
\def\Hom{\mathop{\mathrm{Hom}}\nolimits}
\def\Ext{\mathop{\mathrm{Ext}}\nolimits}
\def\SheafHom{\mathop{\mathcal{H}\!\mathit{om}}\nolimits}
\def\SheafExt{\mathop{\mathcal{E}\!\mathit{xt}}\nolimits}
\def\Sch{\mathit{Sch}}
\def\Mor{\mathop{\mathrm{Mor}}\nolimits}
\def\Ob{\mathop{\mathrm{Ob}}\nolimits}
\def\Sh{\mathop{\mathit{Sh}}\nolimits}
\def\NL{\mathop{N\!L}\nolimits}
\def\CH{\mathop{\mathrm{CH}}\nolimits}
\def\proetale{{pro\text{-}\acute{e}tale}}
\def\etale{{\acute{e}tale}}
\def\QCoh{\mathit{QCoh}}
\def\Ker{\mathop{\mathrm{Ker}}}
\def\Im{\mathop{\mathrm{Im}}}
\def\Coker{\mathop{\mathrm{Coker}}}
\def\Coim{\mathop{\mathrm{Coim}}}

% Boxtimes
%
\DeclareMathSymbol{\boxtimes}{\mathbin}{AMSa}{"02}

%
% Macros for moduli stacks/spaces
%
\def\QCohstack{\mathcal{QC}\!\mathit{oh}}
\def\Cohstack{\mathcal{C}\!\mathit{oh}}
\def\Spacesstack{\mathcal{S}\!\mathit{paces}}
\def\Quotfunctor{\mathrm{Quot}}
\def\Hilbfunctor{\mathrm{Hilb}}
\def\Curvesstack{\mathcal{C}\!\mathit{urves}}
\def\Polarizedstack{\mathcal{P}\!\mathit{olarized}}
\def\Complexesstack{\mathcal{C}\!\mathit{omplexes}}
% \Pic is the operator that assigns to X its picard group, usage \Pic(X)
% \Picardstack_{X/B} denotes the Picard stack of X over B
% \Picardfunctor_{X/B} denotes the Picard functor of X over B
\def\Pic{\mathop{\mathrm{Pic}}\nolimits}
\def\Picardstack{\mathcal{P}\!\mathit{ic}}
\def\Picardfunctor{\mathrm{Pic}}
\def\Deformationcategory{\mathcal{D}\!\mathit{ef}}
\usepackage{cleveref}
\crefname{originallemma}{original lemma}{original lemmas}
\Crefname{originallemma}{Original lemma}{Original lemmas}
\makeatletter
\newlabel{derived-lemma-negative-vanishing}{{05TC}{source}{}{sourcepointer-1}{}}
\newlabel{derived-lemma-negative-vanishing@cref}{{[originallemma][0][]05TC}{[1][0][]source}{}{}{}}
\makeatother
\crefname{originallemma}{original lemma}{original lemmas}
\Crefname{originallemma}{Original lemma}{Original lemmas}
\makeatletter
\newlabel{lemma-pseudo-coherent-independent-representative}{{08CC}{source}{}{sourcepointer-2}{}}
\newlabel{lemma-pseudo-coherent-independent-representative@cref}{{[originallemma][0][]08CC}{[1][0][]source}{}{}{}}
\makeatother
\crefname{originallemma}{original lemma}{original lemmas}
\Crefname{originallemma}{Original lemma}{Original lemmas}
\makeatletter
\newlabel{lemma-strictly-perfect-pullback}{{08H3}{source}{}{sourcepointer-3}{}}
\newlabel{lemma-strictly-perfect-pullback@cref}{{[originallemma][0][]08H3}{[1][0][]source}{}{}{}}
\makeatother
\crefname{originallemma}{original lemma}{original lemmas}
\Crefname{originallemma}{Original lemma}{Original lemmas}
\makeatletter
\newlabel{sites-lemma-morphism-of-sites-covering}{{08H2}{source}{}{sourcepointer-4}{}}
\newlabel{sites-lemma-morphism-of-sites-covering@cref}{{[originallemma][0][]08H2}{[1][0][]source}{}{}{}}
\makeatother
\crefname{originallemma}{original lemma}{original lemmas}
\Crefname{originallemma}{Original lemma}{Original lemmas}
\makeatletter
\newlabel{sites-modules-lemma-localize-morphism-ringed-sites}{{04J0}{source}{}{sourcepointer-5}{}}
\newlabel{sites-modules-lemma-localize-morphism-ringed-sites@cref}{{[originallemma][0][]04J0}{[1][0][]source}{}{}{}}
\makeatother
\crefname{originallemma}{original lemma}{original lemmas}
\Crefname{originallemma}{Original lemma}{Original lemmas}
\makeatletter
\newlabel{sites-modules-lemma-relocalize}{{04IY}{source}{}{sourcepointer-6}{}}
\newlabel{sites-modules-lemma-relocalize@cref}{{[originallemma][0][]04IY}{[1][0][]source}{}{}{}}
\makeatother
\title{The derived pullback of an m-pseudo-coherent object along a morphism of ringed sites is m-pseudo-coherent}
\author{The Stacks Project Authors}
\date{stacks-project:sites-cohomology.tex:local-snapshot-20261011; GNU FDL 1.2 or later}
\begin{document}
\maketitle
\noindent\textit{Editorial note.} All delivered mathematical text is the author's own, copied from the pinned Stacks Project chapter file sites-cohomology.tex; no statement or proof step was rewritten, reordered or removed.\par\medskip
\noindent\textit{Editorial note.} The counted proof is the author's complete argument: the proof is the author's complete argument: the reduction along coverings of both sites, the localisation of the strictly perfect witness, the computation with the distinguished triangle and the vanishing of the pulled-back cone, and the final use of the strictly-perfect pullback lemma are all written out.\par\medskip
\noindent\textit{Editorial note.} The label lemma-pseudo-coherent-pullback names the counted statement, and the proof environment immediately following it in the source is delivered verbatim as its proof.\par\medskip
\noindent\textit{Editorial note.} The statement and proof are copied from the pinned chapter file; the six cited inputs are kept as external pointers with their canonical Stacks tags.\par\medskip
\noindent\textit{Editorial note.} The preamble is the chapter's own retained minimal preamble (sources/stacks-sites-cohomology/retained-preamble.tex), a verbatim excerpt of the native preamble: the Xy-pic setup and the native theorem declarations with the mathematical macros.\par\medskip
\noindent\textit{Allowed standard prerequisites (editorial).}\begin{itemize}
\item m-pseudo-coherent objects of D(O) on a ringed site. Same chapter, the definition of m-pseudo-coherence by the existence of a strictly perfect complex mapping to the object with cohomology isomorphisms above m and a surjection in degree m. Supplies both the hypothesis on E and the property verified for Lf\textasciicircum{}*E.
\item Local character of pseudo-coherence. Same chapter, the independence of the representative and the locality of m-pseudo-coherence under restriction along a covering. Supplies the reduction from an arbitrary U to the members of a covering and to the localisations C/U, D/V.
\item Strictly perfect complexes are pulled back to strictly perfect complexes. Same chapter, the pullback lemma for strictly perfect complexes along a morphism of ringed sites. Supplies the final step: Lf\textasciicircum{}* of a strictly perfect complex is strictly perfect, so it witnesses pseudo-coherence of Lf\textasciicircum{}*E.
\item Vanishing of the left-derived pullback in high degrees. Derived Categories, the vanishing statement for the left derived functor applied to a complex with vanishing cohomology above a bound. Supplies the vanishing H\textasciicircum{}j(Lf\textasciicircum{}*C) = 0 for j >= m, which is what makes the pulled-back morphism a witness.
\end{itemize}\medskip
\ifdefined\hypertarget\hypertarget{sourcepointer-1}{}\fi
\noindent\textit{Source reference.} Original lemma 05TC: sources/stacks-chow/tags.txt, source lines 7476--7476. Prerequisite pointer invoked at proof line(s) 12167; it names a standard fact of a sibling chapter that the counted argument uses, and every step of the counted proof built on it is retained verbatim in the delivered body.\par\medskip
\ifdefined\hypertarget\hypertarget{sourcepointer-2}{}\fi
\noindent\textit{Source reference.} Original lemma 08CC: sources/stacks-sites-cohomology/pinned-original-source.tex, source lines 12066--12066. Prerequisite pointer invoked at proof line(s) 12124; it names a standard fact of the same chapter that the counted argument uses, and every step of the counted proof built on it is retained verbatim in the delivered body.\par\medskip
\ifdefined\hypertarget\hypertarget{sourcepointer-3}{}\fi
\noindent\textit{Source reference.} Original lemma 08H3: sources/stacks-sites-cohomology/pinned-original-source.tex, source lines 11755--11755. Prerequisite pointer invoked at proof line(s) 12173; it names a standard fact of the same chapter that the counted argument uses, and every step of the counted proof built on it is retained verbatim in the delivered body.\par\medskip
\ifdefined\hypertarget\hypertarget{sourcepointer-4}{}\fi
\noindent\textit{Source reference.} Original lemma 08H2: sources/stacks-chow/tags.txt, source lines 10765--10765. Prerequisite pointer invoked at proof line(s) 12121; it names a standard fact of a sibling chapter that the counted argument uses, and every step of the counted proof built on it is retained verbatim in the delivered body.\par\medskip
\ifdefined\hypertarget\hypertarget{sourcepointer-5}{}\fi
\noindent\textit{Source reference.} Original lemma 04J0: sources/stacks-chow/tags.txt, source lines 5921--5921. Prerequisite pointer invoked at proof line(s) 12133; it names a standard fact of a sibling chapter that the counted argument uses, and every step of the counted proof built on it is retained verbatim in the delivered body.\par\medskip
\ifdefined\hypertarget\hypertarget{sourcepointer-6}{}\fi
\noindent\textit{Source reference.} Original lemma 04IY: sources/stacks-chow/tags.txt, source lines 5919--5919. Prerequisite pointer invoked at proof line(s) 12130; it names a standard fact of a sibling chapter that the counted argument uses, and every step of the counted proof built on it is retained verbatim in the delivered body.\par\medskip
\begin{lemma}
\label{lemma-pseudo-coherent-pullback}
Let $(f, f^\sharp) : (\mathcal{C}, \mathcal{O}_\mathcal{C}) \to
(\mathcal{D}, \mathcal{O}_\mathcal{D})$
be a morphism of ringed sites. Let $E$ be an object of
$D(\mathcal{O}_\mathcal{C})$. If $E$ is $m$-pseudo-coherent,
then $Lf^*E$ is $m$-pseudo-coherent.
\end{lemma}

\begin{proof}
Say $f$ is given by the functor $u : \mathcal{D} \to \mathcal{C}$.
Let $U$ be an object of $\mathcal{C}$. By
Sites, Lemma \ref{sites-lemma-morphism-of-sites-covering}
we can find a covering $\{U_i \to U\}$ and for each $i$ a morphism
$U_i \to u(V_i)$ for some object $V_i$ of $\mathcal{D}$.
By Lemma \ref{lemma-pseudo-coherent-independent-representative}
it suffices to show that $Lf^*E|_{U_i}$ is $m$-pseudo-coherent.
To do this it is enough to show that $Lf^*E|_{u(V_i)}$ is
$m$-pseudo-coherent, since $Lf^*E|_{U_i}$ is the restriction
of $Lf^*E|_{u(V_i)}$ to $\mathcal{C}/U_i$ (via
Modules on Sites, Lemma
\ref{sites-modules-lemma-relocalize}).
By the commutative diagram of
Modules on Sites, Lemma
\ref{sites-modules-lemma-localize-morphism-ringed-sites}
it suffices to prove the lemma for the morphism of ringed
sites $(\mathcal{C}/u(V_i), \mathcal{O}_{u(V_i)}) \to
(\mathcal{D}/V_i, \mathcal{O}_{V_i})$.
Thus we may assume $\mathcal{D}$ has a final object $Y$ such that
$X = u(Y)$ is a final object of $\mathcal{C}$.

\medskip\noindent
Let $\{V_i \to Y\}$ be a covering such that for each $i$ there exists
a strictly perfect complex $\mathcal{F}_i^\bullet$ of
$\mathcal{O}_{V_i}$-modules and a morphism
$\alpha_i : \mathcal{F}_i^\bullet \to E|_{V_i}$ of $D(\mathcal{O}_{V_i})$
such that $H^j(\alpha_i)$ is an isomorphism
for $j > m$ and $H^m(\alpha_i)$ is surjective.
Arguing as above it suffices to prove the result for
$(\mathcal{C}/u(V_i), \mathcal{O}_{u(V_i)}) \to
(\mathcal{D}/V_i, \mathcal{O}_{V_i})$. Hence we may assume that
there exists a strictly perfect complex $\mathcal{F}^\bullet$ of
$\mathcal{O}_\mathcal{D}$-modules and a morphism
$\alpha : \mathcal{F}^\bullet \to E$ of $D(\mathcal{O}_\mathcal{D})$
such that $H^j(\alpha)$ is an isomorphism
for $j > m$ and $H^m(\alpha)$ is surjective. In this case, choose
a distinguished triangle
$$
\mathcal{F}^\bullet \to E \to C \to \mathcal{F}^\bullet[1]
$$
The assumption on $\alpha$ means exactly that the cohomology sheaves
$H^j(C)$ are zero for all $j \geq m$. Applying $Lf^*$ we obtain
the distinguished triangle
$$
Lf^*\mathcal{F}^\bullet \to Lf^*E \to Lf^*C \to Lf^*\mathcal{F}^\bullet[1]
$$
By the construction of $Lf^*$ as a left derived functor we see that
$H^j(Lf^*C) = 0$ for $j \geq m$ (by the dual of Derived Categories, Lemma
\ref{derived-lemma-negative-vanishing}). Hence $H^j(Lf^*\alpha)$ is an
isomorphism for $j > m$ and $H^m(Lf^*\alpha)$ is surjective.
On the other hand, since
$\mathcal{F}^\bullet$ is a bounded above complex of flat
$\mathcal{O}_\mathcal{D}$-modules we see that
$Lf^*\mathcal{F}^\bullet = f^*\mathcal{F}^\bullet$.
Applying Lemma \ref{lemma-strictly-perfect-pullback} we conclude.
\end{proof}

\end{document}
